\chapter{Problem Formulation}
\label{ch:problem}

\section{System Model}
An access point receives tasks from ground users and assigns each one to a
satellite. Each task $j$ has an arrival time $a_j$, RAM, disk and image-size
requirements, an execution time $e_j$, a transfer time $t_j$ and a number of
hops. Each satellite $s$ has a CPU capacity, network capacity, latency,
bandwidth, elevation angle, remaining energy and an \emph{operational-until}
time $\omega_s$ (orbital sunset).

\section{Constraints}
A task $j$ placed on satellite $s$ starting at $f_s$ (the time $s$ becomes free)
with duration $d_j = e_j + t_j$ is feasible only if
\begin{align}
  f_s + d_j &\le \omega_s && \text{(operational window)}\\
  \rho_s\, d_j &\le E_s   && \text{(energy budget)}
\end{align}
where $\rho_s$ is the energy rate and $E_s$ the remaining energy. A task is
processed by one satellite and cannot be interrupted.

\section{Objectives}
Given a schedule $\mathcal{A}$ over $N$ tasks:
\begin{align}
  \text{completion ratio}\;\; C &= |\mathcal{A}|/N,\\
  \text{mean response time}\;\; R &= \tfrac{1}{|\mathcal{A}|}\sum_{(j,s)\in\mathcal{A}} (c_j - a_j),\\
  \text{energy}\;\; \mathcal{E} &= \sum_{(j,s)\in\mathcal{A}} \rho_s d_j,
\end{align}
where $c_j$ is the completion time. We maximise $C$ while minimising $R$ and
$\mathcal{E}$.

\section{Simulation Data}
The raw dataset is the task log of one 300-minute LEO simulation (2960 rows, about 92 distinct
servers) produced by the DES simulator of \citet{magliarisi2025orbit}. It is split into 30
consecutive 10-minute windows (``worlds'', 2959 tasks in total, roughly 100 per world). In each
world the arrival times are shifted so that the earliest task arrives at minute 0, keeping the gaps:
the scheduler only depends on the relative arrival pattern. Satellites are built from the server
profiles found in the world's rows. The snapshot time is set to the last arrival, so every task of
the world is available to the decoder. Chapter~\ref{ch:impl} describes the pipeline.
